\scanpage{309}
\mspage{8}
ou encore (moyennant l'équivalence de catégories canonique
\[
 \sigma:(A_*)^{\cdot}\simeq(A^{\cdot})_*\qquad)
\]
un objet $\tfrac12$ simplicial dans la catégorie $A^{\cdot}$ des complexes
de cochaînes, soit
\[
 n\longmapsto {K_{[n]}}^{\cdot}.
\]
On interprète de même les objets des autres catégories mixtes
\deleted{\unsure{$\tfrac12$ simplic. $A^{\cdot*}$}}
\[
 (\unsure{*|R})\quad
 \left\{\begin{array}{llll}
 A_{*\cdot},&A_{\cdot*},&{A^{*\deleted{\unsure{*}}}}_{\cdot},&
 \deleted{\unsure{{A^{\cdot*}}_*}}\\
 A^{*\cdot}_{\deleted{\unsure{*}}},&A^{\cdot*}_{\deleted{\unsure{*}}},&
 {A^{\cdot}}_*,&\deleted{\unsure{{A^*}_{\cdot}}}.
 \end{array}\right.
\]
De façon analogue, les objets $K^{\cdot\cdot}$ de $(A^{\cdot})^{\cdot}$
sont les bicomplexes de cochaînes, les objets $K_{\cdot\cdot}$ de
$(A_{\cdot})_{\cdot}=A_{\cdot\cdot}$ sont les bicomplexes de chaînes,
enfin les objets ${K^{\cdot}}_{\cdot}$ de $(A_{\cdot})^{\cdot}$ ou
${K_{\cdot}}^{\cdot}$ de $(A^{\cdot})_{\cdot}$ sont les bicomplexes mixtes
\inserted{chaînes-cochaînes} (de chaînes pour un indice, de cochaînes pour
un autre). Dans le cas de $K^{\cdot\cdot}$ ou $K_{\cdot\cdot}$ le complexe
simple associé est un complexe de cochaînes resp. de chaînes, tandis que
dans le cas mixte ${K^{\cdot}}_{\cdot}$ ou ${K_{\cdot}}^{\cdot}$, on trouve
un complexe simple avec des objets en degrés positifs ou négatifs à priori
(Il y aurait lieu de fixer une convention de degrés, si pour le degré total
l'op. $\unsure{d}$ est de degré $+1$ ou $-1$, suivant
\margin{310}
